\chapter{Methodology}
\label{chap:methodology}

\section{Research and Development Approach}

The project follows an iterative design-and-evaluate approach. Each phase
produces a component with stable interfaces and tests before the next phase
depends on it. The six phases are: literature and dataset preparation, model
specialization, Adaptive RAG development, evaluation and optimization,
prototype integration, and final reporting. This reduces the risk that a model
experiment becomes inseparable from a particular user interface or storage
implementation.

The unit of comparison in the final evaluation is a saved RAG configuration
applied to a deterministic subset of WixQA. A configuration records route,
classifier, planner, retrieval mode, \textit{top-k}, reranking, citation policy,
generator, fallback, response structure, tone, and prompts. Knowledge-base
embedding remains outside the chat configuration because query vectors must
inherit the active index embedding. This separation allows a configuration to
be reproduced while preventing an invalid query/document vector combination.

\section{Dataset Acquisition and Normalization}

The WixQA download script acquires ExpertWritten, Simulated, Synthetic, and the
knowledge-base corpus. QA rows are normalized to a QAC schema containing:

\begin{itemize}
  \item stable record and source identifiers;
  \item question and reference answer;
  \item supporting context or source-document references;
  \item subset and provenance metadata; and
  \item a complexity label with its provenance.
\end{itemize}

The prepared knowledge corpus is stored as UTF-8 JSONL. Each non-empty row is
treated as one article document with a stable record ID, title, URL, article
type, text, checksum, and source revision. The import pipeline validates the
expected count and duplicate identifiers before creating documents.

The repository audit records 200 ExpertWritten rows, 200 Simulated rows, 6,221
Synthetic rows, and \PreparedCorpusDocuments{} prepared knowledge documents.
Dataset hashes are recorded in the preparation manifest so another environment
can detect a changed source snapshot.

\Cref{tab:wixqa-normalized-examples} illustrates how records with different
origins are represented by the same schema. ExpertWritten answers preserve the
detail of authentic support cases, Simulated answers condense user--expert
dialogues, and Synthetic answers provide larger-scale single-article coverage.
The displayed answers are shortened for readability; the prepared JSONL keeps
the complete normalized text and source identifiers.

\begin{table}[H]
  \centering
  \caption{Representative normalized WixQA records from the prepared dataset.}
  \label{tab:wixqa-normalized-examples}
  \small
  \begin{tabularx}{\textwidth}{
    >{\raggedright\arraybackslash}p{0.13\textwidth}
    >{\raggedright\arraybackslash}p{0.37\textwidth}
    >{\raggedright\arraybackslash}X}
    \toprule
    Subset & Normalized question & Reference-answer excerpt, support, and record ID\\
    \midrule
    ExpertWritten &
    Can I start accepting payments while my Wix Payments account is still
    under verification? &
    Payments can be accepted almost immediately, but identity verification is
    required before full activation. One supporting article is linked.
    \textit{Record:} \Code{wixqa\_expertwritten-00001}.\\
    Simulated &
    How can I keep published changes as a draft? &
    Edit and save the site without selecting Publish. One supporting article
    is linked. \textit{Record:} \Code{wixqa\_simulated-00001}.\\
    Synthetic &
    How can I manage inventory updates after creating purchase orders in the
    Wix dashboard? &
    Inventory updates automatically when received merchandise is recorded in
    Wix Stores POS. One supporting article is linked. \textit{Record:}
    \Code{wixqa\_synthetic-00001}.\\
    \bottomrule
  \end{tabularx}
\end{table}

The corresponding knowledge corpus uses a document rather than a QA schema.
For example, record
\Code{860475a2...09168}
stores the article titled ``Wix Events: About the Event Details and Registration
Form Pages,'' together with its normalized body, public support URL, and article
type. During evaluation, QAC article identifiers are matched to these prepared
documents so retrieval can be assessed against the intended evidence rather
than an unrelated corpus sample.

\section{Four-Class Complexity Dataset}

\subsection{Label Definitions}

\begin{table}[H]
  \centering
  \caption{Operational complexity-label definitions.}
  \label{tab:complexity-labels}
  \begin{tabularx}{\textwidth}{p{0.14\textwidth}p{0.15\textwidth}X}
    \toprule
    Label & Intended route & Evidence requirement\\
    \midrule
    Simple & L1 & Answerable without retrieving organization-specific
    evidence, for example a conversational or general direct request.\\
    Moderate & L2 & Requires one article, one coherent passage, or a single
    retrieval pass. One-document WixQA questions default to this class rather
    than no retrieval.\\
    Complex & L3 & Requires two or three evidence groups, multiple workflow
    steps, decomposition, or synthesis across retrieved passages.\\
    Advanced & L4 & Requires four or more dependencies, conditional or
    exception chains, cross-source investigation, or bounded tool-assisted
    research.\\
    \bottomrule
  \end{tabularx}
\end{table}

The operational distinction is evidence requirement rather than question
length alone. \Cref{tab:four-class-examples} gives representative records from
the prepared classifier artifact. Long questions are abridged only in the
table; their IDs refer to the complete stored records.

\begin{table}[H]
  \centering
  \caption{Representative examples from the prepared four-class dataset.}
  \label{tab:four-class-examples}
  \small
  \begin{tabularx}{\textwidth}{
    >{\raggedright\arraybackslash}p{0.12\textwidth}
    >{\raggedright\arraybackslash}p{0.48\textwidth}
    >{\raggedright\arraybackslash}X}
    \toprule
    Label and route & Representative question & Label rationale and provenance\\
    \midrule
    Simple, L1 &
    ``Define SLA.'' &
    No organization-specific document is needed. The answer is a direct
    workflow-vocabulary definition. Template record
    \Code{template-four-class-00042}; zero evidence documents.\\
    Moderate, L2 &
    ``How can I manage inventory updates after creating purchase orders in the
    Wix dashboard?'' &
    One coherent article directly contains the required fact. Official
    Synthetic record \Code{wixqa\_synthetic-00001}; one evidence document.\\
    Complex, L3 &
    ``I am having trouble finding the arrows and the Change Slide Background
    option in a full-width slideshow on Wix.'' &
    The answer combines slideshow-background instructions with navigation and
    control settings. ExpertWritten record
    \Code{wixqa\_expertwritten-00074}; two evidence documents.\\
    Advanced, L4 &
    Plan an end-to-end Wix Stores workflow combining currency conversion and
    carrier integration, handling a cart-discount exception, and verifying
    multiple-coupon behavior; identify branches, owners, recovery actions, and
    approval order. &
    The request contains a dependency chain, an exception, and final evidence
    across four articles. Abridged template-balancing record
    \Code{template-four-class-02095}; four evidence documents.\\
    \bottomrule
  \end{tabularx}
\end{table}

These examples also clarify an important boundary: a short enterprise lookup
is moderate when it still requires a source, whereas a longer general
definition may remain simple if no private or organization-specific evidence
is needed. Likewise, L4 is not assigned merely because a question is verbose;
it requires an iterative dependency, exception, recovery, or tool-use pattern
beyond L3's bounded multi-passage synthesis.

Labels can be produced through a hybrid silver-label procedure. L1, L2, L3,
and L4 are executed with fixed snapshots. The least-complex successful route
is selected. Answer overlap at or above 0.65 is a deterministic pass, below
0.35 is a deterministic failure, and retrieval routes additionally require a
faithfulness proxy of at least 0.60. Ambiguous outcomes may be resolved by a
registered judge. If execution evidence remains unresolved, structural rules
based on source count and dependency depth provide a fallback label. Every
label records provenance, route outcomes, judge, metrics, cost, latency, and
configuration snapshot.

\subsection{Balancing and Leakage Control}

The current training artifact contains \PreparedDatasetRecords{} records,
balanced at \PreparedDatasetPerClass{} per class. It retains all 200
ExpertWritten and 200 Simulated rows, deterministically samples 200 official
Synthetic rows, and uses 1,800 independently named template examples to fill
class deficits. The template supplement is not represented as official WixQA
Synthetic data.

\begin{figure}[H]
  \centering
  \begin{tikzpicture}
    \begin{axis}[
      ybar,
      width=0.9\textwidth,
      height=7cm,
      ymin=0,
      ymax=700,
      ylabel={Selected records},
      symbolic x coords={Simple,Moderate,Complex,Advanced},
      xtick=data,
      nodes near coords,
      nodes near coords align={vertical},
      bar width=24pt,
      axis line style={draw=aragline},
      tick style={draw=aragline},
      ymajorgrids=true,
      grid style={draw=aragline,dashed},
      every node near coord/.append style={font=\small}
    ]
      \addplot[fill=aragprimary,draw=aragprimary] coordinates {
        (Simple,600) (Moderate,600) (Complex,600) (Advanced,600)
      };
    \end{axis}
  \end{tikzpicture}
  \caption{Balanced four-class classifier dataset after deterministic preparation.}
  \label{fig:dataset-distribution}
\end{figure}


Records are split using source-document groups rather than independent random
rows. Questions that share a source group cannot appear in both training and
validation. With seed 42, the manifest contains \PreparedSourceGroups{} source
groups, \PreparedTrainRecords{} training records, and
\PreparedValidationRecords{} validation records. The validation class counts
are 120 simple, 119 moderate, 119 complex, and 123 advanced. The audit reports
no duplicate IDs, duplicate question groups, conflicting labels, or
train/validation source leakage.

\section{Classifier Training and Routing}

Two model families were trained and evaluated on the same source-grouped
validation split:

\begin{enumerate}
  \item DistilBERT sequence classification with four logits, softmax
  probabilities, confidence, top-two margin, and calibrated validation
  outputs; and
  \item T5-small sequence scoring over the four candidate labels, followed by
  probability normalization.
\end{enumerate}

The common interface preserves \Code{predict(query) -> label} and adds a scored
prediction containing probabilities, confidence, margin, deployment identity,
runtime, and latency. Comparison metrics include accuracy, macro F1,
per-label recall, advanced precision/recall/F1, expected calibration error,
latency, under-routing, over-routing, and route-cost regret.

Adaptive routing maps simple to L1, moderate to L2, complex to L3, and advanced
to L4. A low-confidence prediction may be escalated one route when confidence
is below 0.60 or top-two margin below 0.15. Saved configurations can change
these thresholds. Explicit L1--L4 mode does not use escalation, which is
necessary for controlled route experiments. A selected classifier failure
falls back to the process-default classifier and produces a prominent trace
warning.

\begin{figure}[H]
  \centering
  \resizebox{0.98\textwidth}{!}{%
  \begin{tikzpicture}[node distance=10mm and 12mm]
    \node[aragbox,text width=40mm] (query) {Original user query};
    \node[aragbox,text width=40mm,right=of query] (rewrite) {Conversation-aware standalone query};
    \node[aragbox,text width=40mm,right=of rewrite] (classifier) {Four-class complexity classifier};

    \node[aragdecision,below=17mm of classifier] (route) {Predicted and\\routed label};
    \node[aragbox,text width=31mm,below left=18mm and 58mm of route] (l1) {\textbf{L1 Direct}\\No retrieval};
    \node[aragbox,text width=31mm,below left=18mm and 18mm of route] (l2) {\textbf{L2 Simple RAG}\\One retrieval pass};
    \node[aragbox,text width=31mm,below right=18mm and 18mm of route] (l3) {\textbf{L3 Complex RAG}\\Decompose and aggregate};
    \node[aragbox,text width=31mm,below right=18mm and 58mm of route] (l4) {\textbf{L4 Advanced RAG}\\Bounded planner/tools};
    \node[araglayer,text width=145mm,below=17mm of route,yshift=-27mm] (generate) {Context budgeting, prompt assembly, generation, citation validation, persistence, and trace finalization};

    \draw[aragarrow] (query) -- (rewrite);
    \draw[aragarrow] (rewrite) -- (classifier);
    \draw[aragarrow] (classifier) -- (route);
    \draw[aragarrow] (route) -- node[above left,font=\scriptsize]{simple} (l1);
    \draw[aragarrow] (route) -- node[left,font=\scriptsize]{moderate} (l2);
    \draw[aragarrow] (route) -- node[right,font=\scriptsize]{complex} (l3);
    \draw[aragarrow] (route) -- node[above right,font=\scriptsize]{advanced} (l4);
    \draw[aragarrow] (l1) -- (generate);
    \draw[aragarrow] (l2) -- (generate);
    \draw[aragarrow] (l3) -- (generate);
    \draw[aragarrow] (l4) -- (generate);
  \end{tikzpicture}}
  \caption{Adaptive query routing from L1 to L4. L4 is an extension introduced by this project.}
  \label{fig:adaptive-routes}
\end{figure}


\section{Knowledge and Index Processing}

\subsection{Source Loading and Safety}

The source registry supports TXT, JSON, JSONL, Markdown, PDF, and DOCX uploads,
prepared WixQA corpus records, and public website URLs. Custom \texttt{aratxt},
\texttt{arajson}, and \texttt{aramd} extensions remain disabled until their
schemas are defined. Website loading rejects private-network destinations,
unsafe redirects, oversized responses, unsupported content types, and
timeouts, following SSRF-prevention principles \cite{owaspssrf2026}.

Each document receives source type, URI, MIME type, size, imported timestamp,
title, stable hash, and source-specific metadata. Deduplication occurs before
chunking. Prepared WixQA rows skip file parsing because their text and metadata
have already been normalized.

\subsection{Chunking}

Existing knowledge bases can use fixed-size, sliding-overlap, header-based,
recursive, semantic, parent-child, or structure-aware recursive chunking.
Legacy configurations retain character-based behavior. The WixQA MiniLM
profile uses the exact MiniLM BERT tokenizer and applies the settings in
\Cref{tab:wixqa-chunk-profile}.

\begin{table}[H]
  \centering
  \caption{WixQA MiniLM structure-aware chunking profile.}
  \label{tab:wixqa-chunk-profile}
  \begin{tabularx}{\textwidth}{p{0.34\textwidth}X}
    \toprule
    Parameter & Configured behavior\\
    \midrule
    Parent identity & WixQA article ID; chunks never cross articles.\\
    Target and soft body size & 180 and 210 WordPieces.\\
    Minimum and overlap & 60 and up to 30 WordPieces.\\
    Hard embedding input limit & 240 WordPieces including prefix and special tokens.\\
    Separator order & Heading, subheading, paragraph, list, sentence, token.\\
    Metadata prefix & Title and heading path, capped at 40 WordPieces.\\
    Structural rules & Preserve fitting numbered/bullet lists, merge small
    adjacent chunks, and overlap only inside the same section.\\
    \bottomrule
  \end{tabularx}
\end{table}

Stored chunk text remains clean. The title and section prefix is composed only
for embedding. Diagnostics store article ID, heading path, section ID, body,
prefix, embedding, and overlap token counts, profile ID, forced-split flag, and
below-minimum warning.

\subsection{Embedding and Atomic Indexes}

Knowledge bases select an embedding deployment and whether external processing
is permitted. Ingestion constructs the embedder from that KB configuration.
Remote embedding requires explicit egress permission and confirmation. The
same model, tokenizer limit, dimension, and active index version are applied to
query embeddings.

Indexing is asynchronous and durable. Chunks are embedded and persisted in
batches, a draft index version records the resolved configuration, and
activation occurs only after every batch succeeds. Failure or cancellation
leaves the previous index active. This prevents a query from observing a
partially built or mixed-model index.

\begin{figure}[H]
  \centering
  \resizebox{\textwidth}{!}{%
  \begin{tikzpicture}[node distance=7mm]
    \node[aragbox] (source) {Local files, prepared WixQA corpus, or safe public URL};
    \node[aragbox,right=of source] (load) {Load and normalize document text};
    \node[aragbox,right=of load] (meta) {Extract metadata and content hash};
    \node[aragbox,right=of meta] (dedup) {Deduplicate documents};
    \node[aragbox,below=12mm of dedup] (chunk) {Configured chunking and token diagnostics};
    \node[aragbox,left=of chunk] (embed) {Embedding through active deployment};
    \node[aragbox,left=of embed] (draft) {Persist draft index version in PostgreSQL/pgVector};
    \node[aragdecision,left=of draft] (valid) {All batches\\successful?};
    \node[aragbox,left=of valid] (activate) {Atomically activate new index};

    \draw[aragarrow] (source) -- (load);
    \draw[aragarrow] (load) -- (meta);
    \draw[aragarrow] (meta) -- (dedup);
    \draw[aragarrow] (dedup) -- (chunk);
    \draw[aragarrow] (chunk) -- (embed);
    \draw[aragarrow] (embed) -- (draft);
    \draw[aragarrow] (draft) -- (valid);
    \draw[aragarrow] (valid) -- node[above,font=\scriptsize]{yes} (activate);
    \draw[aragsoftarrow] (valid.south) -- ++(0,-9mm) node[below,font=\scriptsize,align=center]{no: retain previous\\active index};
  \end{tikzpicture}}
  \caption{Versioned knowledge ingestion and index activation workflow.}
  \label{fig:knowledge-pipeline}
\end{figure}


\section{Query-Time Adaptive RAG}

\subsection{Conversation Context and Reformulation}

When enabled, the service loads only completed user--assistant exchanges that
precede the current request. Failed, cancelled, pending, and streaming messages
are excluded. The number of exchanges and character budget are configurable,
subject to server ceilings of six exchanges and 10,000 characters.
Evaluation disables conversation context because each benchmark case is an
independent QA observation.

Referential follow-ups are converted into a standalone query. Deterministic
rewriting is always available. A selected planner can perform the rewrite
through the Model Gateway; invalid output, provider failure, or blocked egress
falls back without failing the answer. The original query is retained for the
answer, while classification, retrieval, reranking, and decomposition use the
standalone query.

\subsection{Retrieval Routes}

L1 calls the generator without knowledge retrieval. L2 performs one configured
retrieval pass. BM25 uses a complete active-index lexical cache; dense search
uses pgVector; hybrid search normalizes and weights both result sets. A
document filter restricts all retrieval to selected document IDs.

L3 creates two to four subqueries, retaining the standalone query as the global
subquery. Each query is retrieved independently. Repeated chunks are
deduplicated by stable identity, then ranked using retrieval score, original
rank, and subquery coverage. Planner decomposition must return a validated JSON
list; otherwise deterministic decomposition is used.

L4 exposes only available, authorized tools: KB hybrid search, chunk/document
inspection, evidence reranking, SSRF-safe public URL retrieval when enabled,
and final-answer completion. The defaults are five planning iterations, eight
tool calls, and 90 seconds. Planner errors are structured observations that may
be recovered. A missing or repeatedly invalid planner falls back to L3.

\subsection{Prompting, Generation, and Citations}

Prompt construction combines system prompt, predefined instructions, response
structure, tone, humor level, bounded conversation history, original question,
standalone query, and labelled evidence. Retrieved contexts receive stable
labels such as \texttt{[S1]}. History and retrieved documents are explicitly
treated as untrusted context rather than system instructions.

The selected generator executes through the Model Gateway. Local extractive
generation provides a deterministic baseline; FLAN-T5 provides optional local
generation; remote or self-hosted models use LiteLLM. A fallback deployment is
permitted only for retryable failure and, during streaming, only before the
first text delta. Citation validation detects missing and unknown labels.
Validation is advisory: it returns a warning rather than triggering a second
paid generation call.

\section{Model Governance Method}

The AI Models catalog separates reusable provider connections from
deployments. Connections contain access path, API base, encrypted credentials,
locality, and health. Deployments contain provider-native model ID,
capabilities, parameters, limits, pricing, budget, health, and enabled status.
Only tested and enabled deployments appear in runtime selectors.

\begin{figure}[H]
  \centering
  \resizebox{\textwidth}{!}{%
  \begin{tikzpicture}[node distance=9mm and 13mm]
    \node[araglayer,text width=50mm] (consumer) {RAG, knowledge worker, evaluation, and analytics consumers};
    \node[aragbox,text width=50mm,below=of consumer] (policy) {Deployment resolution, capability validation, budget and egress policy};
    \node[araglayer,text width=50mm,below=of policy] (gateway) {Provider-neutral Model Gateway};
    \node[aragbox,text width=42mm,below left=of gateway] (builtin) {LocalBuiltinAdapter\\Extractive, FLAN-T5, hash, MiniLM, reranker, classifiers};
    \node[aragbox,text width=42mm,below right=of gateway] (lite) {LiteLLMAdapter\\Generate, stream, embed, rerank, judge, planner};
    \node[aragbox,text width=27mm,below=of lite,xshift=-34mm] (exp) {OpenRouter\\experimentation};
    \node[aragbox,text width=27mm,below=of lite] (direct) {OpenAI/Gemini\\direct};
    \node[aragbox,text width=27mm,below=of lite,xshift=34mm] (self) {Ollama/vLLM\\self-hosted};

    \draw[aragarrow] (consumer) -- (policy);
    \draw[aragarrow] (policy) -- (gateway);
    \draw[aragarrow] (gateway) -- (builtin);
    \draw[aragarrow] (gateway) -- (lite);
    \draw[aragarrow] (lite) -- (exp);
    \draw[aragarrow] (lite) -- (direct);
    \draw[aragarrow] (lite) -- (self);
  \end{tikzpicture}}
  \caption{Unified AI Model Gateway and runtime access paths.}
  \label{fig:model-gateway}
\end{figure}


Credentials entered through the UI are encrypted using a versioned AES-GCM
envelope. API responses expose only credential readiness, never the value.
Each model attempt records purpose, deployment, connection, requested and
resolved model, latency, tokens, estimated cost, status, fallback index, and
related conversation, KB, or evaluation identifiers. Embeddings never fall
back to a different model because that would violate vector compatibility.

\section{Evaluation Design}

An experiment selects one or more saved RAG configurations, one or more WixQA
subsets, a knowledge base, a judge deployment, a case limit, and deterministic
sampling seed. Each configuration--dataset cell runs independently. There is
no hidden static-L2 baseline; Adaptive and L2 are compared by selecting two
saved configurations.

\begin{table}[H]
  \centering
  \caption{Evaluation dimensions and representative metrics.}
  \label{tab:evaluation-metrics}
  \begin{tabularx}{\textwidth}{p{0.24\textwidth}X}
    \toprule
    Dimension & Metrics\\
    \midrule
    Routing & Accuracy, macro F1, confusion matrix, per-label recall,
    under-routing, over-routing, route distribution, route-cost regret.\\
    Retrieval & Hit@\textit{k}, MRR, precision@\textit{k}, recall@\textit{k},
    NDCG@\textit{k}, retrieval coverage, context relevancy, context recall.\\
    Generation & Token F1, BLEU, ROUGE-1, ROUGE-2, answer overlap,
    factuality, context adherence, answer relevancy, grading note.\\
    Operations & End-to-end and stage latency, tokens, estimated cost, model
    calls, retrieved contexts, failures, cancellations, and fallback use.\\
    Human feedback & Per-answer-version thumbs up/down and optional comment,
    linked to configuration, KB, user, and trace.\\
    \bottomrule
  \end{tabularx}
\end{table}

Judge calls use structured output and one malformed-JSON repair attempt. Runs
remain failed rather than silently complete when required metrics are absent.
RAGXplain diagnosis is a separate bounded job over selected cases. This avoids
unexpected judge costs and allows normal evaluation results to remain valid if
diagnosis fails.

\begin{figure}[H]
  \centering
  \resizebox{\textwidth}{!}{%
  \begin{tikzpicture}[node distance=8mm]
    \node[aragbox] (select) {Select KB, WixQA subset, RAG configuration, and judge};
    \node[aragbox,right=of select] (queue) {Queue durable evaluation experiment};
    \node[aragbox,right=of queue] (execute) {Run each case through configured L1--L4 pipeline};
    \node[aragbox,below=12mm of execute] (metrics) {Compute retrieval, lexical, judge, latency, token, and cost metrics};
    \node[aragbox,left=of metrics] (persist) {Persist run, cases, snapshots, traces, and leaderboard};
    \node[aragbox,left=of persist] (diagnose) {Run bounded RAGXplain diagnosis};
    \node[aragbox,below=of diagnose] (viewer) {Embedded executive insights and configuration actions};

    \draw[aragarrow] (select) -- (queue);
    \draw[aragarrow] (queue) -- (execute);
    \draw[aragarrow] (execute) -- (metrics);
    \draw[aragarrow] (metrics) -- (persist);
    \draw[aragarrow] (persist) -- (diagnose);
    \draw[aragarrow] (diagnose) -- (viewer);
  \end{tikzpicture}}
  \caption{Configuration-based WixQA evaluation and RAGXplain diagnosis workflow.}
  \label{fig:evaluation-workflow}
\end{figure}


\section{Validity, Reproducibility, and Ethics}

Internal validity is supported by immutable configuration snapshots,
deterministic sampling, grouped splits, artifact hashes, route traces, and
recorded model identities. Construct validity remains limited because lexical
overlap and LLM judges approximate answer quality. External validity is limited
to Wix business support content and the configurations actually tested.
Conclusion validity remains limited by the 20-case compatible sample reported
in \Cref{chap:implementation-results}; larger controlled samples are required
for statistical claims.

The repository records the Git commit, Python version, model deployment, KB
index, prompts, and evaluation settings. Full traces improve reproducibility
but may contain prompts and retrieved enterprise content. They are retained for
a bounded period, compressed, size-limited, and sanitized to exclude secrets,
authorization headers, encrypted credentials, and raw numerical vectors.
Remote processing is blocked unless enabled for the KB, providing an explicit
data-egress decision.
